% Independent draft supplement: --main supplementary.tex --no-build-script.
% Standard article layout accommodates full-width tables without the journal's
% production-only title banner or two-column table constraints.
\documentclass[11pt,a4paper]{article}
\usepackage[margin=25mm]{geometry}
\usepackage{amsmath,amssymb,graphicx,array,booktabs,tabularx,siunitx,longtable}
\usepackage[numbers,sort&compress]{natbib}
\usepackage[hidelinks]{hyperref}
\sisetup{uncertainty-mode=separate}
\renewcommand{\thefigure}{S\arabic{figure}}
\renewcommand{\thetable}{S\arabic{table}}
\renewcommand{\theequation}{S\arabic{equation}}

\title{Supplementary data: Manuscript title}
\author{First Author and Second Author\\
  \small Department, Institution, City, Country\\
  \small \texttt{first.author@example.org}}
\date{}

\begin{document}
\maketitle
Replace this instructional overview with a brief description of the supplied
supplementary material. This document is independent of the main manuscript.
All example values and graphics are synthetic.

\section{Supplementary methods}
Describe supporting methods and derivations. Figures, tables and equations
use independent S numbering. Equation~\eqref{eq:si-residual} demonstrates a
multi-line expression with one number.
\begin{equation}
\begin{split}
 r_i&=y_i-\widehat y_i,\\
 E&=\frac{1}{n}\sum_{i=1}^{n}r_i^2.
\end{split}\label{eq:si-residual}
\end{equation}
This is a notation example, not a scientific finding.

\section{Supplementary figures}
Figure~\ref{fig:si-slot} demonstrates a supporting figure. Replace the box
and legend with the study's own graphic, units and uncertainty definitions.
\begin{figure}[htbp]
\centering
\fbox{\parbox[c][48mm][c]{0.80\linewidth}{\centering
Supplementary figure slot\\[8pt]Replace with an original PDF or PNG}}
\caption{\textbf{Supporting graphic.} Replace this legend with a description
of the actual figure. The box is a layout example and represents no data.}
\label{fig:si-slot}
\end{figure}
\clearpage

\section{Supplementary tables}
Table~\ref{tab:si-numeric} illustrates aligned numerical columns.
Table~\ref{tab:si-long} repeats column headings when a supporting inventory
extends over multiple pages.
\begin{table}[htbp]
\centering\small
\caption{Synthetic entries used to illustrate numerical alignment.}
\label{tab:si-numeric}
\begin{tabularx}{\linewidth}{@{}>{\raggedright\arraybackslash}X
 S[table-format=2.0] S[table-format=1.3(3)]@{}}
\toprule
Setting & {Count} & {Estimate $\pm$ SD}\\
\midrule
Example A & 24 & 0.320 +- 0.015\\
Example B & 36 & 0.415 +- 0.021\\
\bottomrule
\end{tabularx}
\par\smallskip\footnotesize\raggedright
All values are invented. Replace them with verified study results and define
the actual measure of uncertainty.
\end{table}
\clearpage

\small
\begin{longtable}{@{}p{0.13\linewidth}p{0.27\linewidth}p{0.50\linewidth}@{}}
\caption{Synthetic supporting inventory to demonstrate continuation.}
\label{tab:si-long}\\
\toprule Item & Setting & Explanation\\\midrule
\endfirsthead
\multicolumn{3}{@{}l}{Table~\thetable\ (continued)}\\[3pt]
\toprule Item & Setting & Explanation\\\midrule
\endhead
\midrule\multicolumn{3}{r@{}}{Continued on next page}\\
\endfoot
\bottomrule
\endlastfoot
S01 & Identifier & Record the stable identifier and origin of each supporting input.\\
S02 & Version & Record the exact release used rather than a changing label.\\
S03 & Provenance & Explain how the input was obtained and whether access is restricted.\\
S04 & Units & Define the measurement units used in the source.\\
S05 & Inclusion & State the rule used to select source records.\\
S06 & Exclusion & State the rule used to exclude source records.\\
S07 & Partition & Explain separation of development and evaluation material.\\
S08 & Mapping & Preserve links between source records and derived observations.\\
S09 & Processing & Describe preprocessing and retain the original source.\\
S10 & Scaling & Define the reference quantities used for normalisation.\\
S11 & Parameters & Record analysis settings and the reason for each choice.\\
S12 & Initialisation & Record initial conditions or random seeds where relevant.\\
S13 & Repetition & Define an independent repetition without duplicating dependent data.\\
S14 & Counts & State the actual number of observations in each comparison.\\
S15 & Uncertainty & Define the uncertainty measure and its calculation.\\
S16 & Interval & State the interpretation of any reported interval.\\
S17 & Comparison & Identify the reference condition for each comparison.\\
S18 & Controls & Describe the controls that support interpretation.\\
S19 & Validation & Describe independent validation where it was performed.\\
S20 & Sensitivity & State settings varied in the study's sensitivity analysis.\\
S21 & Software & Record exact software versions and relevant options.\\
S22 & Environment & Record computational details that affect reproducibility.\\
S23 & Precision & Retain the precision justified by the source measurements.\\
S24 & Rounding & Explain display rounding without changing source values.\\
S25 & Aggregation & Define how observations are combined and weighted.\\
S26 & Dependence & Identify observations sharing the same subject or source.\\
S27 & Labels & Check figure labels at the intended final size.\\
S28 & Figure source & Retain plotting code and editable original graphics.\\
S29 & Table source & Link displayed values to verified analysis output.\\
S30 & Symbols & Define mathematical variables, operators and dimensions.\\
S31 & Assumptions & State the conditions required by each derivation.\\
S32 & References & Check supplementary citations independently of the main paper.\\
S33 & Availability & Provide an accurate access route and any restrictions.\\
S34 & Limitations & State limitations supported by the study design and evidence.\\
\end{longtable}
\normalsize

\section{Supplementary references}
This public citation tests the independent bibliography\cite{wilkinson2016fair};
it is not evidence for any synthetic example above.
\bibliographystyle{unsrt}
\bibliography{references}
\end{document}
